%% 第三章 强化学习算法设计

\chapter{强化学习算法设计}
\chaptermark{强化学习算法设计}

本章阐述用于训练自驱动颗粒迁移策略的强化学习算法设计。
首先介绍强化学习的基本概念和算法分类，在此基础上引入Soft Actor-Critic (SAC) 算法，
重点阐述其熵正则化机制、网络架构与参数更新方法；
随后给出面向点对点迁移任务的状态空间、动作空间与奖励函数的具体设计；
最后介绍训练流程与超参数配置方案。

\section{强化学习基本框架}

强化学习是一类通过智能体与环境交互来学习最优决策策略的机器学习方法。
在每一个离散时间步$t$，智能体观测环境状态$s_t$，依据策略$\pi(a_t | s_t)$
选择动作$a_t$并执行，环境据此返回新状态$s_{t+1}$和即时奖励信号$r_t$。
智能体的学习目标是最大化累积期望回报

\begin{equation}
    J(\pi) = \mathbb{E}_{\pi}\left[\sum_{t=0}^{\infty} \gamma^t r_t\right],
    \label{eq:rl_objective}
\end{equation}

其中$\gamma \in (0, 1]$为折扣因子，用以平衡近期奖励与远期奖励的相对重要性。

按照优化目标的不同，现代强化学习算法可大致分为基于值函数的方法和
基于策略优化的方法两类。基于值函数的算法（如DQN）
通过学习最优动作值函数$Q^*(s,a)$的近似$Q_\theta(s,a)$间接确定策略，
其优化过程与具体策略无关；基于策略梯度的算法（如PPO）
则将策略直接参数化为$\pi_\theta(a|s)$，通过对目标$J(\pi_\theta)$做梯度上升
来优化策略参数。将值函数方法与策略优化方法相结合的Actor-Critic架构
能够同时利用二者优势，在连续控制任务中取得了优异的性能。

在本研究的场景中，Actor-Critic架构的优势尤为明显。
自推进体需要在连续二维空间中选择推进方向和速度大小，
这是一个典型的高维连续控制问题，
基于值函数的方法（如DQN）需要对动作空间进行离散化处理，
不仅导致动作粒度的精度损失，
还会引发维度灾难——若将二维方向均匀划分为$N$个离散角度、
速度大小划分为$M$个等级，
则动作空间大小为$N \times M$，
而Actor-Critic架构通过策略网络直接输出连续动作参数，
从根本上规避了离散化的精度与效率问题。
此外，双涡流流场的非定常特性意味着环境动态性是时变的，
离线策略学习方法（如SAC所采用的）
能够有效复用历史经验数据，
避免因环境动态变化导致的旧数据失效问题。


在连续控制领域，基于最大熵框架的Soft Actor-Critic算法已成为处理高维连续动作空间的主流方法。
Zhao等验证了SAC算法的最大熵准则在连续控制任务中的优势地位，
解析了其通过最大化策略熵来平衡探索与利用的内在机制\cite{zhaoSACPathPlanning2023}。
Yu等进一步将SAC网络架构与人工势场法相结合进行监督学习，
使得导航计算时间与路径冗余障碍大幅缩减\cite{yuSoftActorcriticCombining2025}。
针对目标导向任务中极度缺乏正向奖励反馈的难题，
已有研究将事后经验回放与SAC算法深度结合，大幅提升了稀疏反馈环境下的采样效率\cite{zhaoSACPathPlanning2023}。
为克服传统人工势场容易导致智能体陷入局部极小值死锁的缺陷，
Bi等引入了可变方向势场机制，允许网络牺牲局部安全性进行大幅度探索以实现全局逃逸\cite{biVDPFPathPlanning2025}。
在前沿控制器设计中，Mohammad等将控制势垒函数作为SAC运行时的底层安全过滤器，
通过动态拦截和重组高危控制动作，实现了极端工况下的鲁棒物理防护\cite{mohammadCBFSAC2025}。
上述算法演进为本文在非定常双涡流环境下构建自推进体能耗优化迁移策略奠定了理论与工程基础。
\section{Soft Actor-Critic 算法}

\subsection{熵正则化强化学习}

Soft Actor-Critic（SAC）是一种以非策略梯度方式优化随机策略的
离线Actor-Critic算法。
区别于标准强化学习仅最大化累积奖励，SAC的核心特征是在优化目标中
引入策略的熵正则化项。

熵是信息论中度量随机变量不确定性的基本概念。设随机变量$x$服从分布$P$，
则其熵$H$定义为

\begin{equation}
    H(P) = \underset{x \sim P}{\mathbb{E}}\left[-\log P(x)\right].
    \label{eq:entropy_def}
\end{equation}

在熵正则化强化学习中，智能体在每个时间步获得的奖励附加与策略熵相关的正则项，
优化目标由此变为

\begin{equation}
    \pi^{*} = \arg\max_{\pi} \underset{\tau \sim \pi}{\mathbb{E}}\left[
        \sum_{t=0}^{\infty} \gamma^{t}
        \left(R(s_t, a_t, s_{t+1}) + \alpha H(\pi(\cdot | s_t))\right)
    \right],
    \label{eq:entropy_rl_obj}
\end{equation}

其中$\alpha > 0$为温度系数，控制熵项相对于奖励项的权重。

引入熵正则化具有两方面的优势。其一，促进探索——
增大熵的权重等价于鼓励智能体输出更为多样的动作，
可有效防止策略过早收敛至局部最优解。
其二，增强鲁棒性——最优策略倾向于在多个相近的优质动作之间维持一定的概率分布，
而非退化为单一确定性动作。在湍流这类具有非定常扰动的环境中，
这一特性对策略的泛化能力尤为重要。

在本研究中，熵正则化的上述两重优势与非定常双涡流环境的特性高度契合。
首先，双涡流场的时间周期性振荡意味着流场结构处于持续演化之中，
智能体如果过早收敛至确定性策略，
将难以适应不同振荡相位下的流场配置——
熵正则化通过强制策略保持一定的随机性，
使智能体在训练过程中持续探索各种可能的流场状态与动作组合，
从而覆盖更广泛的状态-动作空间。
其次，在流场环境与迁移任务的双重随机性下
（流场的非定常演化与起终点位置的随机选取），
最优策略本身可能具有多模态特征——
在不同的流场条件下，
可能存在多条能耗相近的迁移路径。
确定性策略会将所有概率质量集中于单一动作，
而随机策略能够自然地维持多条备选路径的决策分布，
在遭遇未知流场扰动时展现出更强的适应性与鲁棒性。
因此，SAC算法的最大熵框架
与本研究面临的非定常流场能耗优化问题在本质上具有高度的适配性。


\subsection{网络架构}

SAC算法同时维护五个神经网络：

(1) 策略网络（Actor）$\pi_\theta(a_t | s_t)$：以状态$s_t$为输入，
输出动作分布的参数。动作从高斯分布中采样后经tanh函数压缩至
动作空间的合法范围内。

(2) 两个Q函数网络（Critic）$Q_{\theta_1}(s_t, a_t)$和$Q_{\theta_2}(s_t, a_t)$：
分别估计在状态$s_t$下执行动作$a_t$的期望累积回报。
采用双Q网络设计可以减轻Q值估计中的系统性高估问题：
更新时以两个Q网络中较小的值作为目标值。

(3) 两个目标Q网络$Q_{\bar{\theta}_1}$和$Q_{\bar{\theta}_2}$：
其参数通过指数移动平均从对应Q网络软更新而来，
用于计算时序差分目标值以提高训练稳定性。

\subsection{损失函数与参数更新}

对于Q网络的训练，SAC采用均方Bellman误差（MSBE）作为损失函数。
以经验回放池$\mathcal{D}$中采样的转移数据为依据，
每个Q函数$Q_{\theta_i}$的损失函数为

\begin{equation}
    L(\theta_i, \mathcal{D}) = \mathbb{E}_{(s, a, r, s', d) \sim \mathcal{D}}\left[
        \left(Q_{\theta_i}(s, a) - y(r, s', d)\right)^2
    \right],
    \label{eq:q_loss}
\end{equation}

目标值$y$的计算同时融入了熵项和双Q网络的最小值操作

\begin{equation}
    y(r, s', d) = r + \gamma(1 - d)\left(
        \min_{j=1,2} Q_{\bar{\theta}_j}(s', \tilde{a}') -
        \alpha \log \pi_\theta(\tilde{a}' | s')
    \right),
    \label{eq:target_value}
\end{equation}

其中$\tilde{a}' \sim \pi_\theta(\cdot | s')$表示从当前策略中重新采样的动作，
$d$为回合终止标志。

策略网络$\pi_\theta$的优化方向为最大化Q值估计与熵项的组合，
其梯度上升所使用的梯度为

\begin{equation}
    \nabla_\theta \frac{1}{|\mathcal{B}|} \sum_{s \in \mathcal{B}}\left(
        \min_{i=1,2} Q_{\theta_i}(s, \tilde{a}_\theta(s)) -
        \alpha \log \pi_\theta(\tilde{a}_\theta(s) | s)
    \right),
    \label{eq:pi_loss}
\end{equation}

其中$\mathcal{B}$为从经验回放池$\mathcal{D}$中随机抽取的一个批次样本。
本文采用固定的熵正则化系数$\alpha = 0.1$，不随训练过程自适应调节。
关于SAC算法的更多细节可参见文献\cite{haarnojaSoftActorcriticOffpolicy2018}。

\section{奖励函数设计}

奖励函数决定了智能体从环境获得的反馈信号，是连接任务目标与策略优化的关键。
本文将迁移任务的完成与能耗的最小化统一在标量奖励信号中，
设计了由三个分量构成的复合奖励函数

\begin{equation}
    r(t) = r_s(t) + r_e(t) + r_h(t),
    \label{eq:reward_total}
\end{equation}

其中$r_s$、$r_e$和$r_h$分别为状态奖励、能量奖励和时间奖励，
以下逐一给出各分量的定义。

状态奖励$r_s$基于颗粒的当前位置与目标点之间的几何关系。
若颗粒迁移出流场腔体的有效区域，则给予$-10$的惩罚；
若颗粒相对于上一时间步更加靠近目标点，即
$\|\mathbf{x}_{\text{agent}}^{t+1} - \mathbf{x}_{\text{goal}}\|_2 <
\|\mathbf{x}_{\text{agent}}^{t} - \mathbf{x}_{\text{goal}}\|_2$，
则获得一个基础正向奖励$e_{\text{basic}}$；
其余情况下奖励为零。因此$r_s$可表示为

\begin{equation}
    r_s = \begin{cases}
        -10, & \text{颗粒迁移出腔体范围},\\
        e_{\text{basic}}, &
        \|\mathbf{x}_{\text{agent}}^{t+1} - \mathbf{x}_{\text{goal}}\|_2 <
        \|\mathbf{x}_{\text{agent}}^{t} - \mathbf{x}_{\text{goal}}\|_2,\\
        0, & \text{其他}.
    \end{cases}
    \label{eq:state_reward}
\end{equation}

能量奖励$r_e$基于颗粒在迁移过程中因自主驱动而产生的能量消耗。
设颗粒的推进速度$\mathbf{u}_{\text{propel}}$与当地背景流场速度
$\mathbf{u}_{\text{fluid}}$之间的夹角为$\alpha$。
当二者方向一致（即$0^{\circ} \leq \alpha \leq 90^{\circ}$）时，
颗粒获得正奖励$e_{\text{basic}} + (e_{\max} - e)$，
其中$e = \frac{1}{2}\|\mathbf{u}_{\text{propel}}\|^2$表示当前步的推进动能，
$e_{\max} = \frac{1}{2}(\|\mathbf{u}_{\text{propel}}\|_{\max})^2$为最大允许推进动能，
$e_{\text{basic}} = e_{\max}$。
这意味着在顺流条件下，颗粒输出的推进速度越低，获得的能量奖励越高。
反之，当颗粒逆流行进（$90^{\circ} < \alpha \leq 180^{\circ}$）时，
将受到惩罚$-(e_{\text{basic}} + e)$，推进速度越高则惩罚越重。
综上，$r_e$可表示为

\begin{equation}
    r_e = \begin{cases}
        e_{\text{basic}} + (e_{\max} - e), &
        0^{\circ} \leq \alpha \leq 90^{\circ},\\
        -(e_{\text{basic}} + e), &
        90^{\circ} < \alpha \leq 180^{\circ}.
    \end{cases}
    \label{eq:energy_reward}
\end{equation}

时间奖励$r_h$基于颗粒完成任务所需的时间消耗。
若颗粒在最大允许时间$t_{\max}$内未能到达目标点，则给予$-5$的惩罚；
若颗粒已进入目标点的邻域（即$\|\mathbf{x}_{\text{agent}}^{t} -
\mathbf{x}_{\text{goal}}\|_2 < \delta_0$，其中$\delta_0$为一个小阈值），
则认为颗粒成功到达终点，给予正奖励$\varepsilon(t_{\max} - t)$，
该奖励与剩余可用时间成正比，颗粒到达得越早奖励越高；
其余情况奖励为零。因此$r_h$可表示为

\begin{equation}
    r_h = \begin{cases}
        -5, & t \geq t_{\max},\\
        \varepsilon(t_{\max} - t), &
        \|\mathbf{x}_{\text{agent}}^{t} - \mathbf{x}_{\text{goal}}\|_2 < \delta_0,\\
        0, & \text{其他}.
    \end{cases}
    \label{eq:time_reward}
\end{equation}

上述奖励设计体现了如下的权衡逻辑：在单个时间步内，
当颗粒既向目标靠近、又顺流行进时，由于$e_{\text{basic}} + (e_{\max} - e) > e_{\text{basic}}$，
减少能量消耗是当前步的主要优化目标，接近终点为次要目标。
但在整个回合尺度上，由于时间奖励的累积期望大于状态与能量奖励之和，
即$\mathbb{E}[\sum r_h] > \mathbb{E}[\sum (r_s + r_e)]$，
尽快抵达目标点始终是首要目标，能耗最小化次之。
这种双层目标设计使智能体在高效完成迁移任务的同时，
自主学习利用背景湍流降低自身推进能耗。

这种“单步节能优先、回合到达优先”的设计策略具有明确的行为引导意图。
在训练初期，智能体尚未掌握任何流场知识，
如果能耗惩罚过强，
智能体可能倾向于“原地不动”以减少能耗，
导致任务完全无法完成。
通过设置$r_h$的时间递减奖励和$r_s$的距离递减奖励，
即使智能体的推进策略极其低效，
只要最终到达目标点，
也能获得正向累积奖励——
这为早期探索提供了必要的正向引导。
随着训练推进，智能体逐渐学会利用流场结构，
$r_e$的节能引导开始发挥主导作用：
在保证到达的前提下，
智能体不断优化推进方向与流场方向的夹角，
使能耗持续降低。

从物理视角看，三个奖励分量分别对应了无人机滑翔任务中的三个核心约束：
$r_s$对应空间约束（必须到达指定目标），
$r_e$对应能量约束（以最小能耗完成迁移），
$r_h$对应时间约束（在合理时间内完成）。
这三个约束之间存在天然的竞争关系——
过于追求时间最优可能导致大量“对抗”流场的能耗，
过于追求能耗最优可能导致迂回路径过长而超出时间限制。
本文的三元复合奖励结构通过权重参数实现了三者的显式折中，
使智能体能够在训练中自动寻找帕累托最优解。

\section{训练流程与超参数配置}

\subsection{训练流程}

SAC算法的训练流程如算法\ref{alg:training}所述。
训练开始前，随机初始化策略网络参数$\theta$和两个Q网络参数$\theta_1, \theta_2$，
并将目标网络参数同步为$\bar{\theta}_1 \leftarrow \theta_1, \bar{\theta}_2 \leftarrow \theta_2$。
在每个环境交互步中，智能体观测当前状态$s$，
依据策略网络采样动作$a \sim \pi_\theta(\cdot | s)$并执行，
获得下一状态$s'$、奖励$r$和终止标志$d$，
将转移元组$(s, a, r, s', d)$存入经验回放池$\mathcal{D}$。
当回放池中样本数量达到预设的最小阈值后，
每隔固定步数从中随机采样一批数据用于网络参数的更新。

具体地，每次更新时从$\mathcal{D}$中抽取批大小为$N_{\text{batch}}$的样本$\mathcal{B}$，
按式(\ref{eq:target_value})计算目标Q值，通过梯度下降最小化MSBE损失以更新两个Q网络，
再通过单步梯度上升优化策略网络，最后以指数移动平均的方式软更新目标网络。
回合终止后环境重置，进入下一回合。

\begin{algorithm}[htbp]
    \caption{SAC训练流程}
    \label{alg:training}
    \begin{algorithmic}
        \Statex \textbf{初始化：}
        \State 随机初始化策略网络参数 $\theta$ 及两个Q网络参数 $\theta_1, \theta_2$
        \State 初始化空的经验回放池 $\mathcal{D}$
        \State 同步目标网络：$\bar{\theta}_1 \leftarrow \theta_1,\; \bar{\theta}_2 \leftarrow \theta_2$
        \Statex
        \For{$t = 0, 1, 2, \dots$}
            \State 观测状态 $s$，采样动作 $a \sim \pi_\theta(\cdot | s)$ 并执行
            \State 接收下一状态 $s'$、奖励 $r$ 及终止标志 $d$
            \State 将 $(s, a, r, s', d)$ 存入 $\mathcal{D}$
            \If{$d$ 为真}
                \State 重置环境
                \State \textbf{continue}
            \EndIf
            \If{回放池样本充足}
                \For{更新次数}
                    \State 从 $\mathcal{D}$ 中采样批数据 $\mathcal{B}$
                    \State 按式(\ref{eq:target_value})计算目标值 $y(r, s', d)$
                    \State 梯度下降更新Q网络参数 $\theta_i$
                    \State 梯度上升更新策略网络参数 $\theta$
                    \State 软更新目标网络：$\bar{\theta}_i \leftarrow \tau \bar{\theta}_i + (1-\tau) \theta_i$
                \EndFor
            \EndIf
        \EndFor
    \end{algorithmic}
\end{algorithm}

\subsection{超参数配置}

神经网络采用全连接前馈架构，策略网络和Q网络均包含两个隐藏层，
每层256个神经元，激活函数使用tanh。
经验回放池容量为$10^6$条经验，批采样大小为256。
本文采用的SAC算法使用固定的熵正则化系数$\alpha = 0.1$，
不随训练过程自适应调整。

关于温度系数$\alpha$的取值，本文选择固定为0.1而非采用可调温度机制（如自动调整$\alpha$以维持目标熵值），
主要基于以下考量：其一，在双涡流场中，探索与利用的平衡需求在训练全程中相对稳定——
流场的周期性振荡保证了环境始终具有足够的随机性以维持探索需求；
其二，固定温度系数简化了超参数调优流程，
避免了自动温度调节可能引入的训练不稳定因素。
关于网络架构，每层256个神经元的设定在表达能力与计算效率之间取得了折中：
更大的网络可能带来更高的表示精度，
但也会增加训练时间和过拟合风险；
更小的网络则可能无法充分捕捉流场状态与最优动作之间的复杂非线性映射关系。主要超参数汇总于表\ref{tab:hyperparams}。

\begin{table}[!h]
    \caption{本文中SAC算法使用的主要超参数}
    \label{tab:hyperparams}
    \centering
    \begin{tabular}{l|c}
        \hline
        参数 & 取值 \\ \hline
        学习率 & $3 \times 10^{-4}$ \\
        经验回放池容量 & $10^6$ \\
        网络更新间隔步长 & 100 \\
        批采样大小 & 256 \\
        折扣因子 $\gamma$ & 0.99 \\
        熵正则化系数 $\alpha$ & 0.1 \\
        \hline
    \end{tabular}
\end{table}

\section{本章小结}

本章系统阐述了面向自驱动颗粒能耗优化迁移任务的强化学习算法设计。
首先介绍了强化学习的基本框架和算法分类，在此基础上引入SAC算法，
重点阐述了熵正则化机制的原理与优势：通过最大化策略的熵来促进探索并增强鲁棒性。
随后详细给出了包含策略网络和双Q网络的五网络架构，
以及基于均方Bellman误差和策略梯度上升的参数更新方法，
采用固定熵正则化系数$\alpha = 0.1$与tanh激活函数。
本文进一步分析了熵正则化与非定常湍流环境的内在适配关系：
流场的时间非定常性与任务的随机性要求策略具备持续探索能力和多模态决策分布，
这正是最大熵框架的核心优势所在。

奖励函数采用状态奖励$r_s$、能量奖励$r_e$和时间奖励$r_h$三元结构：
$r_s$基于颗粒与目标点的相对位置关系，$r_e$基于推进速度与背景流方向的夹角，
$r_h$基于任务完成的时间消耗。该设计在单步尺度上偏向能耗最小化，
在回合尺度上保证任务优先完成，实现了迁移效率与能量节省的有效折中。
本文进一步揭示了三个奖励分量与无人机滑翔任务中
空间约束、能量约束和时间约束之间的对应关系，
阐明了三元复合结构如何通过显式权重实现多目标帕累托优化。

从算法与物理问题的适配性角度看，
SAC算法在连续动作空间中的高效探索能力
与自推进体在非定常流场中需要平滑、连续调整推进方向的物理需求高度匹配；
熵正则化提供的鲁棒性则为策略在流场振荡条件下的泛化提供了保障；
而基于局部观测的状态空间设计使得算法框架具有良好的可扩展性，
可便捷地迁移至更复杂的流场环境。
最后给出了完整的训练流程伪代码和超参数配置表，
为后续仿真实验提供了算法实现基础。

关于SAC算法中双Q网络设计的动机，值得在此深入讨论。
在标准Q-learning中，Q值的高估问题
是一个已被充分认识的现象——
由于Q网络的更新目标中包含了最大化操作，
估计误差会系统性地偏向高估方向。
在非定常流场环境中，由于环境动态性的时变特征，
状态-动作值函数的估计面临更大的不确定性，
Q值高估的风险更为突出。
SAC通过取两个独立Q网络的最小值来构建目标值，
这一Clipped Double-Q技巧有效地将估计偏差从高估
校正为轻微的保守估计，
从而避免了智能体对某些高回报
但不稳健的动作的过度依赖。
这一设计对于本文场景尤为重要：
流场涡旋结构可能使某些瞬时动作看起来极具吸引力
（如被涡旋快速推向目标方向），
但如果没有双Q网络的保守估计机制，
智能体可能过度信任这类非鲁棒的高回报经验，
导致策略在实际部署时表现不稳定。

